\thesisAppendix{Г}{Фрагменты ключевого кода}

В настоящем приложении приведены шесть фрагментов исходного кода клиентских приложений системы <<ВКурсе>> --- по два на каждую платформу (iOS, Android и веб-составляющая). Фрагменты подобраны в соответствии с архитектурными решениями, рассмотренными в главе~3, и иллюстрируют те места кодовой базы, где сосредоточена нетривиальная инженерная логика, не уместившаяся в основной текст работы. Для каждого листинга в предшествующем абзаце указан путь к файлу в репозитории, диапазон строк и ссылка на соответствующий пункт главы~3.

Двойная визуальная адаптация iOS-приложения под iOS~26+ Liquid Glass, обоснованная в п.~\ref{subsubsec:impl-ios-liquid-glass}, реализована через систему парных ресурсов \texttt{AccentColor}: для каждой из двенадцати акцентных палитр в каталоге ассетов заведены парные \texttt{colorset}-ы (\texttt{AccentBlue} для iOS~26+ и \texttt{AccentBlueLegacy} для iOS~17.6--18 и т.\,д.). Выбор между ними сосредоточен в одном вычисляемом свойстве \texttt{assetName}, ветви которого срабатывают через \texttt{\#available(iOS 26.0, *)}; одно и то же имя ресурса переиспользуется и в SwiftUI-ветке (\texttt{color}), и в UIKit-ветке (\texttt{uiColor}), что исключает рассинхронизацию темы между двумя ветками рендеринга и масштабируется на двенадцать акцентов без дублирования кода. Файл \texttt{Shared/Core/Common/Design/AccentColor.swift}, строки 10--58 (листинг~\ref{lst:ios-accent-color}).

\begin{lstlisting}[caption={Парные colorset'ы \texttt{AccentColor} для двойной адаптации под iOS~26+ Liquid Glass},label=lst:ios-accent-color]
enum AccentColor: String, CaseIterable, Identifiable {
    case blue = "Blue"
    case cyan = "Cyan"
    case green = "Green"
    case pink = "Pink"
    case orange = "Orange"
    case mint = "Mint"
    case indigo = "Indigo"
    case red = "Red"
    case yellow = "Yellow"
    case purple = "Purple"
    case chartreuse = "Chartreuse"
    case loving = "Loving"

    // ... unrelated: applySystemTint(for:) for UIKit mirroring ...

    var id: String { rawValue }

    var assetName: String {
        if #available(iOS 26.0, *) {
            "Accent\(rawValue)"
        } else {
            "Accent\(rawValue)Legacy"
        }
    }

    var uiColor: UIColor? { UIColor(named: assetName) }
    var color: Color { Color(assetName) }
}
\end{lstlisting}

Дедупликация расписания при смене подписок, описанная в п.~\ref{subsubsec:impl-ios-data}, реализована методом \texttt{deleteLessons(in:for:)} репозитория \texttt{LessonRepository}. Реализация учитывает, что один и тот же \texttt{LessonDTODataModel} может одновременно <<принадлежать>> нескольким сущностям --- учебной группе, преподавателю и аудитории, --- и поэтому прямое удаление по предикату ключа уничтожило бы записи, ещё нужные другим активным подпискам. Алгоритм сначала собирает множество оставшихся активных ключей из \texttt{ScheduleMetadataModel} (исключая ключ удаляемой подписки) и удаляет занятие только в том случае, если оно не совпадает ни с одним из них. Файл \texttt{Shared/Core/Data/Repositories/LessonRepository.swift}, строки 122--153 (листинг~\ref{lst:ios-lesson-repo-dedup}).

\begin{lstlisting}[caption={Дедупликация уроков при удалении подписки в \texttt{LessonRepository}},label=lst:ios-lesson-repo-dedup]
private func deleteLessons(
    in context: ModelContext,
    for keyToDelete: LessonEntityKey
) {
    let allMetadata =
        (try? context.fetch(FetchDescriptor<ScheduleMetadataModel>())) ?? []
    let activeKeys =
        allMetadata
        .filter {
            $0.entityTypeRaw != keyToDelete.type.rawValue
                || $0.entityID != keyToDelete.entityID
        }
        .map {
            LessonEntityKey(
                type: ScheduleType(rawValue: $0.entityTypeRaw)!,
                entityID: $0.entityID
            )
        }

    let predicate = makePredicate(for: keyToDelete)
    let descriptor = FetchDescriptor<LessonDTODataModel>(
        predicate: predicate
    )
    let lessons = (try? context.fetch(descriptor)) ?? []

    for lesson in lessons {
        let isUsedElsewhere = activeKeys.contains {
            lesson.matches(key: $0)
        }
        if !isUsedElsewhere { context.delete(lesson) }
    }
}
\end{lstlisting}

Типобезопасная декларативная навигация Android-приложения, описанная в п.~\ref{subsubsec:impl-android-app-model}, задана единым файлом \texttt{Screens.kt}: иерархия \texttt{object} и \texttt{data class}, аннотированных \texttt{@Serializable} из библиотеки \texttt{kotlinx.serialization}, перечисляет все точки назначения \texttt{NavController}. Параметры маршрута передаются как поля типа, а не как строковые ключи вида \texttt{"schedule/\{id\}?type=\{type\}"}: компилятор гарантирует тип и состав аргументов, а сериализация при переходе по маршруту выполняется автоматически. Файл \texttt{app/\allowbreak src/\allowbreak main/\allowbreak java/\allowbreak app/\allowbreak vcourse/\allowbreak vcourse/\allowbreak core/\allowbreak navigation/\allowbreak Screens.kt}, строки 12--108 (листинг~\ref{lst:android-screens}).

\begin{lstlisting}[caption={Типизированная декларативная навигация Compose через \texttt{@Serializable}},label=lst:android-screens]
package app.vcourse.vcourse.core.navigation

import app.vcourse.vcourse.core.common.enums.FavoritesCategory
import app.vcourse.vcourse.core.common.enums.ScheduleType
import kotlinx.serialization.Serializable

// MARK: ONBOARDING

@Serializable
object ScreenWelcomeOnboarding

@Serializable
object ScreenUniversitySelectionOnboarding

@Serializable
object ScreenRoleSelectionOnboarding

// MARK: HOME

@Serializable
object ScreenHome

// ... unrelated ...

@Serializable
data class ScreenClassroomList(val buildingJson: String?, val addToRecent: Boolean = false)

// MARK: LESSON

@Serializable
data class ScreenLesson(val routeLessonJson: String, val type: ScheduleType)

@Serializable
data class ScreenNote(val routeNoteJson: String)

// ... unrelated ...

// MARK: SCHEDULE

@Serializable
data class ScreenSchedule(val id: Int, val type: ScheduleType, val addToRecent: Boolean = false)

// MARK: DATA CLASSES

@Serializable
data class RouteBuilding(val id: String, val name: String)

@Serializable
data class RouteLesson(
    val slot: Int,
    val dateHashCode: Int,
    val entityId: Int,
    val scheduleType: ScheduleType
)

@Serializable
data class RouteNote(val id: String, val title: String, val text: String?)
\end{lstlisting}

Точка входа Android-приложения совмещает три архитектурных аспекта в единой композиционной функции \texttt{AppEntryPoint}: типизированную навигацию Compose (см. листинг~\ref{lst:android-screens}), реактивную смену темы и акцентного цвета (Material~3 Expressive, см.~п.~\ref{subsubsec:impl-android-ui}) и горячее переключение языка через \texttt{LanguageManager.createLocalizedContext(...)}, который вызывает системный \texttt{createConfigurationContext} и не требует перезапуска Activity (Android~10+). Activity дополнительно объявлена в манифесте с атрибутом \texttt{android:\allowbreak configChanges=\allowbreak "uiMode\allowbreak |orientation\allowbreak |screenSize\allowbreak |screenLayout\allowbreak |smallestScreenSize"}, что предотвращает её пересоздание при перечисленных изменениях конфигурации и сохраняет состояние графа навигации (\texttt{rememberNavController()}) и моделей (\texttt{rememberAppModel()}) без надстройки \texttt{ViewModel}. Файл \texttt{app/\allowbreak src/\allowbreak main/\allowbreak java/\allowbreak app/\allowbreak vcourse/\allowbreak vcourse/\allowbreak MainActivity.kt}, строки 68--94 (листинг~\ref{lst:android-app-entry}).

\begin{lstlisting}[caption={Точка входа Android-приложения с реактивной темой и горячим переключением языка},label=lst:android-app-entry]
@Composable
private fun AppEntryPoint(state: InitialAppState) {
    val initialThemeMode = state.themeMode ?: AppDefaults.THEME_MODE
    val initialThemeColor = state.themeColor ?: AppDefaults.THEME_COLOR
    val initialLanguage = state.language ?: AppDefaults.LANGUAGE

    val currentThemeMode by ThemeModeManager.themeModeFlow.collectAsState(initialThemeMode)
    val currentThemeColor by ThemeColorManager.themeColorFlow.collectAsState(initialThemeColor)
    val currentLanguage by LanguageManager.languageFlow.collectAsState(initialLanguage)

    val localizedContext = remember(currentLanguage) {
        LanguageManager.createLocalizedContext(ContextHolder.appContext, currentLanguage)
    }

    val appModel = rememberAppModel()

    CompositionLocalProvider(
        LocalLocalizedContext provides localizedContext,
        LocalAppModel provides appModel
    ) {
        AppTheme(currentThemeColor, currentThemeMode) {
            Surface(color = MaterialTheme.colorScheme.surface) {
                AppNavGraph(rememberNavController())
            }
        }
    }
}
\end{lstlisting}

BFF-прокси веб-составляющей, обоснованный в п.~\ref{subsubsec:impl-web-bff}, реализован одним catch-all-обработчиком \texttt{[...path]/route.ts}. Обработчик инкапсулирует всю интеграцию с приватным backend-API <<ВКурсе>>: дозаписывает обязательные заголовки \texttt{Authorization: Bearer ...} (токен \texttt{VC\_AUTH\_TOKEN} живёт только в серверном рантайме и не утекает в браузер) и \texttt{X-API-Version}, переносит query-параметры клиента, опирается на встроенный Data Cache Next.js через \texttt{next:~\{~revalidate:~3600~\}}. Тривиальный по объёму, но архитектурно ключевой код: именно он превращает фронтенд в защищённый BFF и снимает с клиентских компонентов знание о схеме внешнего API. Файл \texttt{src/app/api/vcourse/[...path]/route.ts} (листинг~\ref{lst:web-bff-proxy}).

\begin{lstlisting}[caption={BFF-прокси к серверной части с инъекцией заголовков аутентификации и версии API},label=lst:web-bff-proxy]
import { NextRequest, NextResponse } from "next/server";

export async function GET(
  request: NextRequest,
  { params }: { params: Promise<{ path: string[] }> },
) {
  const { path } = await params;
  const targetPath = path.join("/");

  const url = new URL(`${process.env.VC_API_URL}/${targetPath}`);

  request.nextUrl.searchParams.forEach((value, key) => {
    url.searchParams.set(key, value);
  });

  if (targetPath === "universities") {
    url.searchParams.set("available", "true");
  }

  const response = await fetch(url.toString(), {
    next: { revalidate: 3600 },
    headers: {
      Authorization: `Bearer ${process.env.VC_AUTH_TOKEN}`,
      "X-API-Version": process.env.VC_API_VERSION || "2",
      Accept: "application/json",
    },
  });

  if (!response.ok) {
    return NextResponse.json(
      { error: "Failed to fetch" },
      { status: response.status },
    );
  }

  const data = await response.json();
  return NextResponse.json(data);
}
\end{lstlisting}

Стек \emph{undo}/\emph{redo} редактора расписания в сервисе администрирования, описанный в п.~\ref{subsubsec:impl-web-admin}, реализован хуком \texttt{useScheduleHistory}: ограниченная глубиной \texttt{HISTORY\_LIMIT~=~100} двойная очередь past/future хранится в \texttt{useRef}, а дискриминированное объединение \texttt{ScheduleHistoryEntryBody} с четырьмя видами действий (\texttt{atoms-append}, \texttt{atoms-remove}, \texttt{atoms-replace}, \texttt{entity-select}) обеспечивает exhaustive-проверку TypeScript в \texttt{switch}-разборе записи истории. Параметр \texttt{expectedId} в функции \texttt{undo} позволяет всплывающему уведомлению <<отменить>> откатывать строго ту запись, которой оно соответствует, даже если поверх неё уже были выполнены другие операции. Файлы \texttt{src/.../schedule/context/history-types.ts}, строки 6--16 (тип записи), и \texttt{src/.../schedule/context/use-schedule-history.ts}, строки 11--110 (хук); приведены совместно в листинге~\ref{lst:web-undo-redo}.

\begin{lstlisting}[caption={Ограниченный стек undo/redo редактора расписания с дискриминированным объединением действий},label=lst:web-undo-redo]
export type ScheduleHistoryEntryBody =
  | { kind: "atoms-append"; atoms: LessonAtom[] }
  | { kind: "atoms-remove"; atoms: LessonAtom[] }
  | { kind: "atoms-replace"; before: LessonAtom[]; after: LessonAtom[] }
  | {
      kind: "entity-select";
      before: SelectedEntity | null;
      after: SelectedEntity | null;
    };

// ----- use-schedule-history.ts -----

const HISTORY_LIMIT = 100;

// ... unrelated (ScheduleHistoryApply / ScheduleHistoryState types) ...

let historyCounter = 0;

function pushBounded(
  arr: ScheduleHistoryEntry[],
  entry: ScheduleHistoryEntry,
): ScheduleHistoryEntry[] {
  const next = [...arr, entry];
  return next.length > HISTORY_LIMIT
    ? next.slice(next.length - HISTORY_LIMIT)
    : next;
}

// ... unrelated (hook declaration, past/future refs, syncFlags) ...

const recordHistory = useCallback(
  (body: ScheduleHistoryEntryBody): number => {
    const id = ++historyCounter;
    const entry = { ...body, id } as ScheduleHistoryEntry;
    pastRef.current = pushBounded(pastRef.current, entry);
    futureRef.current = [];
    syncFlags();
    return id;
  },
  [syncFlags],
);

const undo = useCallback(
  (expectedId?: number) => {
    const past = pastRef.current;
    if (past.length === 0) return;
    const entry = past[past.length - 1];
    if (expectedId !== undefined && entry.id !== expectedId) return;
    switch (entry.kind) {
      case "atoms-append":
        applyAtomsRemove(entry.atoms.map((a) => a.id));
        break;
      case "atoms-remove":
        applyAtomsAppend(entry.atoms);
        break;
      case "atoms-replace":
        applyAtomsReplace(
          entry.after.map((a) => a.id),
          entry.before,
        );
        break;
      case "entity-select":
        applySelectEntity(entry.before);
        break;
    }
    pastRef.current = past.slice(0, -1);
    futureRef.current = pushBounded(futureRef.current, entry);
    syncFlags();
  },
  // ... unrelated deps ...
);

// ... unrelated (symmetric redo) ...
\end{lstlisting}
